\chapter{网络通信与机器人控制}
\setchapterheader{第四章\quad 网络通信与机器人控制}
\label{chap:network-control}
\begin{intro}
\begin{itemize}
\item 分清 IP、网卡、子网、路由与端口的职责。
\item 了解 ping、DDS 和手柄输入的检查方法。
\item 用示例程序认识观测、策略、动作和控制模式。
\end{itemize}
\end{intro}

机器人控制程序需要同时处理网络、手柄和模型。网络传输机器人的状态和指令，手柄提供运动输入，模型则根据观测计算动作。本章介绍这几部分怎样连接起来，以及怎样检查它们的工作情况。

\section{IP 地址与网络接口}
一台笔记本可以同时有 Wi-Fi 和有线网卡。IP 配置附着在具体接口上，因此先用 \code{ip -br addr} 找实际有线网卡名，再把指定地址配置到它上面。手册中的 \code{enp3s0} 只是例子。

\code{192.168.55.101/24} 表示这个 IPv4 地址带有 24 位网络前缀，对应掩码 \code{255.255.255.0}。本部署中，笔记本的 \code{.101} 和运控板的 \code{.102} 位于同一 \code{192.168.55.0/24} 子网；两个设备不能同时占用同一个地址。

\begin{command}
ip -br addr
ip route get 192.168.55.102
\end{command}
第二条显示内核为目标地址选择的路由、网络接口和源地址。配置完成后，应看到流量走实际有线网卡，源地址为 \code{192.168.55.101}。

\begin{concept}{网关、DNS 与本机直连}
默认网关用于没有更具体路由时转发流量；DNS 把域名解析成 IP。直接访问同子网的数字 IP 通常不需要 DNS，也不需要给机器人直连网卡配置默认网关。下载依赖所需的互联网连接，与机器人有线子网可以分别检查。
\end{concept}

\section{网络通信与 DDS}
\code{ping} 向目标发送 ICMP 请求，并显示响应时间与丢包情况。配置好 IP 后，我们先用它检查连通性，再检查 DDS 的硬件状态回调。

DDS 配置文件选择 UDP、关闭组播，并把运控板地址作为静态 Peer。DDS 负责程序间数据交换；IP 指定主机，端口用于区分传输端点。DDS 的通信参数保存在 \code{config/dds.xml} 中，排查时可以先检查这个文件。

\begin{longtable}{@{}P{0.22\linewidth}P{0.67\linewidth}@{}}
\toprule
\textbf{观察到的结果} & \textbf{检查内容} \\
\midrule
能下载 Python 包 & 用于下载依赖的互联网连接正常。\\
\code{ping} 有回复 & ICMP 可达；还需检查 DDS 配置、接口选择和硬件状态回调。\\
\code{jstest} 数值变化 & 接收器和本机输入设备已有数据。\\
模型加载成功日志 & ONNX 模型相关初始化已进行；还需验证硬件状态和手柄。\\
\bottomrule
\end{longtable}

在 \code{process()} 中，\code{latest_hw_status is None} 或手柄数据为 \code{None} 都会让函数提前返回。所以“程序还在跑，按键却没反应”时，应分别定位网络状态与本机输入。

\section{策略模型与控制指令}
策略是把观测映射为动作的函数；ONNX 是本项目保存模型的一种格式，ONNX Runtime 负责执行它。部署是加载已有策略并进行推理；训练则是产生或更新策略参数的过程。本指南处理前者。

示例程序把运动指令、IMU 与关节状态、上一帧动作组成观测；模型给出 21 维动作。动作再经缩放，与默认关节位置组合成目标位置：
\[
q_i^{\mathrm{des}} = q_i^{\mathrm{default}} + s_i a_i .
\]
其中 $a_i$ 是模型输出，$s_i$ 是对应关节的动作比例。程序根据这个公式生成目标位置，并把位置、刚度、阻尼等参数一起提交给 SDK。

\begin{remark}{更换策略模型}
第一次部署使用 SDK 配套的模型与配置。以后更换模型时，需要逐项对应 21 维动作、$71\times5$ 维观测、关节顺序、默认角度、动作比例和控制增益。示例程序使用固定关节列表，因此模型的关节排列也要与它一致。
\end{remark}

\section{控制频率与模式切换}
主循环目标周期为 $0.002$ s，即名义 500 Hz；在正常策略循环中，每 10 次更新一次模型动作，对应名义 50 Hz。两次策略更新之间，程序继续向 SDK 提交相应关节指令。这是“控制频率”与“策略更新频率”的区别。

\[
f_{\mathrm{policy}} \approx \frac{f_{\mathrm{loop}}}{\mathrm{decimation}}
=\frac{500}{10}=50\ \mathrm{Hz}.
\]
实际周期还受到程序耗时和系统调度的影响。打印大量日志、设备读写阻塞或后台任务较多时，循环周期会变长。因此运行前要接通电源、开启性能模式，并减少后台负载。

阻尼、零位插值、站立准备和策略控制，是程序组织控制行为的不同模式。后文第 8 步介绍如何通过手柄依次切换这些模式。
